\preludeandfugue{The Last House}
\contributor{Ella Low}

% BIOG: 
% 
% Ella Low is a Romanian-born British writer of dark fiction and
% poetry. She has a background in Philosophy, which probably explains
% her tendency to question everything and her fascination with people,
% death and the stranger aspects of being alive. She is the author of
% Too Many Shades of Dark and Ten Dark Stories and lives in
% Worcestershire, where she continues to write about things most people
% would rather not think about.

\prelude

The house was the first of its kind to be built on the vast lands
owned by Mr George Coombe. As a modern man, with progressive ideas and
a love of symmetry and order, he wanted to build a house that
reflected those qualities. His namesake, George~I, was ruling Britain at the
time, and George took that as an omen as well as an incentive; his
architectural progeny would become part of history. His legacy would
be there for entire generations of Coombes to come.

It took the best part of six years for the actual building to be
completed. This was not a rushed job. The house had to be strong and
sound, and built to last for centuries. Each brick was carefully dug
from the clay of the estate and hand-fired by skilled builders. Then
talented plasterers, stonemasons and painters worked diligently to
transform the shell into a masterpiece. And a masterpiece it was.

From the grand hallway, the beautiful staircase cantilevered
effortlessly towards the first floor, and the high-ceilinged rooms
were tastefully arranged. Spacious and luminous, with large and
mathematically precise six-by-six sash windows, the one-storey house
stood grand. Robust yet refined, it was a fine example of understated
elegance, symmetry and tranquillity.
 
\fugue

I remember the day my very first brick was put in place. I was built
to impress: careful proportions, solid foundations, an imposing
staircase. I was built with strong bones. I was built to last. Then
came the people. I was to contain generations of noise, chaos, and
heartbreak. Humans can be infuriating creatures, showing precious
little regard for the good things on this land. Witnessing so many
births, weddings and deaths led me to regard humanity as a peculiar
sort of infestation.

At first, the inhabitants wore bulky clothes and they were quieter,
but the tragedies seemed never-ending. I have watched tiny coffins
leave through my front door. I have learned the weight of mourning
crepe upon my mirrors. With an astonishing lack of personal hygiene
and common sense, it is little wonder that many of them never reached
old age.

I have seen brother killing brother, young marrying old, nannies
smothering babes, and for what\=to claim ownership over me. As though
I could ever be owned. They’ve tried to claim me in other, smaller
ways over the centuries, too: notches in the doorframes measuring
their offspring every year; hiding beautiful walls behind their
hideous, vain portraits; redecorating with garish wallpaper every
generation or so; changing light fittings and modifying entire rooms
with complete disregard for proportion or beauty.

As time went by, the clothes became skimpier, the voices louder and
more entitled, the marriages less about land and lineage, and more
about love, and more of them lived long enough to grow old.

Yet, decade after decade, generation after generation, century after
century, I’ve been able to detect a pattern – beyond their foibles and
follies. Humans, as a species, are very simple and vulnerable
creatures. Once I’d realised that, I became less consumed by their
fleeting lifetimes.

I was built to last, and their desperate life achievements washed past me.

No matter how much they’ve tried, they couldn’t destroy my sturdy
stone heart. `This house has good bones,' they said. I accepted the
nail holes, the scuff marks, the outrageous colour schemes, the
children’s pencil drawings behind the wallpaper. I decided to hold
them, learn about them. I learned the rhythm of their footsteps and
bore witness to their milestones. I listened to prayers whispered into
pillows and arguments hurled across dinner tables. I protected them
from winter and sheltered them from storms.

Then… the flash: the Sun the humans invented, brighter than noon.

That morning, everything changed. In that instant, my family left
behind shadows on my bleached walls. They didn’t get to see much of
what followed: the grand staircase turned into a pile of rubble. Great
oak doors scorched to ash. My beautiful windows smashed into a million
pieces. Some walls gave way and crumbled immediately; some did so
slowly, piece by piece, brick by brick, losing a bit of hope year
after year.

The rain enters freely now where the west wall used to stand.

The cold came early, and somehow decided not to leave. It’s been
dreadfully still since. Everything is ashen and drab.

There are no more sparrows to nest in my eaves, and it’s only the moss
that is growing again, clinging to me.

There are nights when the silence is too quiet, when the cold stars
are visible through the ribs of my broken roof.

The shadows on the wall have started to fade, but I hope that wall
will never collapse.

This place has good bones, my last people said.

And so did they.


